\section{The Base: From the Integers}
\label{sec:integers}

If reality is one experiential field, how do we describe its structure without smuggling in
anything arbitrary? This is where care is needed, because a single confusion can derail the
whole picture.

\paragraph{Experience is the substance, the integers are the description.}
The foundation of reality is experience, the field of vibes. It is not the integers. Integers
do not feel anything. But to describe the \emph{structure} of the field precisely, we need a
language, and we want a language that adds nothing of its own. Every theory that begins with
space, or particles, or a list of tuned constants has quietly assumed all of those, and they
could all have been otherwise. The least arbitrary mathematical structure there is, the one
with nothing to choose, is the whole numbers. So the integers are the right \emph{skeleton}
to draw the field with, precisely because they import no assumptions. The slogan ``the
universe from the integers'' is a statement about the cleanest description, not about the
substance. Experience is the flesh and the life. The integers are the bones.

\paragraph{The ladder.}
From the integers, a short climb produces the substrate, each rung adding exactly one idea.

\begin{enumerate}
\item \textbf{The integers.} The bedrock of the description. No settings, no choices.
\item \textbf{Their symmetries.} The ways the integers can be rearranged without changing
their structure form a clean object (the modular group), built from two simple moves. We did
not invent it. We asked the integers for their symmetries, and this is what came back.
\item \textbf{From symmetries to mirrors.} Widen the idea to reflections, as in a
kaleidoscope, where a few mirrors turn one scrap of color into an endless pattern. The
mathematics of kaleidoscopes is the theory of Coxeter groups, and the only data they need is
the angles between the mirrors, which are again whole numbers.
\item \textbf{A curved space, tiled.} Let the kaleidoscope run in hyperbolic space, the
saddle-shaped geometry that keeps opening outward, and the reflected copies of one tile fill
the space with a crystal.
\item \textbf{The beautiful crystals.} Choosing the integers selects a specific crystal, the
heptagrid of sevens, the pentagrid of fives, the three-dimensional dodecagrid. These are among
the most beautiful objects in mathematics, and any of them can serve as the fabric of
spacetime.
\end{enumerate}

\paragraph{The golden thread.}
Something quietly astonishing happens along the way. The golden ratio, near 1.618, the number
of sunflowers and seashells, keeps appearing in the foundation, and from independent
directions: it is the most even way to scatter points, it governs how the crystal grows ring
by ring, and it is the central, straightest path through the crystal's addressing scheme.
Three separate roads arrive at the same number. When a constant turns up that many times
without being put there on purpose, it is a sign one is looking at something genuinely
ordered and whole, not a patchwork.

\paragraph{The picture in one line.}
The integers, given their symmetries, become a kaleidoscope. The kaleidoscope, set loose in
curved space, becomes a crystal. The crystal, grown by a simple rule and dressed in feeling,
becomes a universe. Reality, on this view, is a mesh of vibes, and experiential field,
modelled by arithmetic, automatically generated into highly geometric structure.
The foundation has minimal arbitrary things in it, just the base axioms assuming experience and a minimal
representation of their qualities (currently focusing on the ternary encoding), and no patterns require randomness.
It is the whole numbers and the elegant structure they carry, all the way up, with experience as the thing
that structure is the structure of.
